\documentclass[11pt]{article}

\usepackage[margin=1in]{geometry}
\usepackage{booktabs}
\usepackage{longtable}
\usepackage{graphicx}
\usepackage{amsmath}
\usepackage[font=small,labelfont=bf]{caption}
\usepackage{float}
\usepackage{hyperref}

\title{CS4372 Assignment 1\\\large Bike Sharing --- SGDRegressor and OLS}
\author{N. Ohayon Rozanes}
\date{\today}

\begin{document}
\maketitle

The data set chosen for this assignment is the Bike Sharing data set
from UC Irvine's machine learning repository. First, the variables
are inspected and the preprocessing process is described. Then a
linear regression is fit with SGDRegressor using SciKit-Learn, then
an OLS and penalized OLS using statsmodels. The goal of the analysis
is to determine mobility patterns in a city as a function of the time
of day, the point in the year, and weather.

AI use disclosure: AI was only used in the production of this report in
the formatting of the tables and figures present in the report.

\section{Feature distributions}
\label{sec:features}

\begin{figure}[H]
  \centering
  \includegraphics[width=\textwidth]{figures/distributions.pdf}
  \caption{Marginal distribution of every numeric feature in the raw UCI
  Bike Sharing (hourly) data. Bins are computed per facet.}
  \label{fig:distributions}
\end{figure}

\begin{figure}[H]
  \centering
  \includegraphics[width=\textwidth]{figures/timeseries.pdf}
  \caption{Hourly ridership (\texttt{cnt}) over the full observation window;
  the band is the intraday spread around each day's mean.}
  \label{fig:timeseries}
\end{figure}

The dataset essentially has two varieties of predictor variables,
categorical temporal variables and continuous weather related
variables, with the exception of \textit{weathersit}, which is weather
related but ordinal. The dataset also includes a timeseries
element, reporting day by day counts. Of course, since time
dependence breaks independence assumptions, trying to include the day
as a regression variable would be problematic. Further, since we are
more interested in movement patterns depending on the qualities of a
day rather than the performence of the company, the specific date
does not capture any additional information that other variables do
not report on.

\textit{workingday} is defined as if it is not a weekend or a
holiday, that is, it is perfectly colinear with the \textit{holiday}
and \textit{weekday} column, and can therefore be safely dropped from
the analysis.

Since we are not interested in the time trend of the data,
\textit{yr} is not qualitatively interesting in the analysis, but as
Figure~\ref{fig:timeseries} shows, there is a noticible phase shift
from 2011 to 2012.
Omitting \textit{yr} would cause omitted variable bias, as such, it
is included in the model.

\textit{atemp} is the ``actual'' temperature that day, and is a
function of the measured temperature, humidity, and windspeed, among
other factors, this is supported by its high VIF. It is therefore
omitted from the model, since it does not capture any additional
information on the bike rides as a function of temperature. A further
analysis into the relationship between the weather related variables
is presented in the PCA section.

The choice between month and season, which are obviously
correlated to each other, is largely arbitrary. They present the same
information at different coarsness. For this analysis, it was found
that the additional values of the month did not carry any additional
information when measured with AIC and BIC.

Lastly, a data entry error seems to report \textit{weatherit}, which
is an ordinal measure of how cloudy the day is, as four seperate
values, but the page itself only lists three. This is likely an error
since very few days are actually reported as having a value of four.
Level four was merged into the level three.

Tables~\ref{tab:vif} and~\ref{tab:cond} demonstrate the positive
effects of the drops on the
data. Making it much more suitable for a regression analysis.
Categorical variables were ran through one-hot encoding and standard
scaling. One hot encoding was chosen in place of ordinal since its
plausible that the relationship between the hour and the ridership is
not necessarily linear, this is confirmed after fitting in
Figure~\ref{fig:hour-effects}.

\begin{table}[htbp]
\centering
\caption{Variance inflation factors for the non-dummy predictors, before and after the collinearity drops. \texttt{---} marks a predictor removed from the final matrix.}\label{tab:vif}
\begin{tabular}{lrr}
\toprule
predictor & before drops & after drops \\
\midrule
yr & 1.02 & 1.02 \\
mnth & 3.41 & --- \\
holiday & 1e+15 & 1.09 \\
workingday & 1e+15 & --- \\
temp & 48.78 & 3.10 \\
atemp & 45.29 & --- \\
hum & 1.84 & 1.82 \\
windspeed & 1.23 & 1.16 \\
\bottomrule
\end{tabular}
\end{table}

\begin{table}[htbp]
\centering
\caption{Collinearity of the design matrix as a whole, before and after the drops.}\label{tab:cond}
\begin{tabular}{lrr}
\toprule
diagnostic & before drops & after drops \\
\midrule
max VIF over dummy levels & 1e+15 & 4.23 \\
condition number & 3.33e+15 & 38.4 \\
\bottomrule
\end{tabular}
\end{table}


\begin{figure}[H]
  \centering
  \includegraphics[width=0.85\textwidth]{figures/corr_heatmap.pdf}
  \caption{Correlation matrix of the standardised design matrix.}
  \label{fig:corr}
\end{figure}

\section{Principal component analysis: weather block}
\label{sec:pca}

\begin{table}[htbp]
\centering
\caption{PCA on the standardised weather block (\texttt{temp}, \texttt{atemp}, \texttt{hum}, \texttt{windspeed}): eigenvalues and explained variance.}\label{tab:pca-variance}
\begin{tabular}{lrrr}
\toprule
component & eigenvalue & proportion of var & cumulative \\
\midrule
PC1 & 1.9966 & 0.4991 & 0.4991 \\
PC2 & 1.2904 & 0.3226 & 0.8217 \\
PC3 & 0.7017 & 0.1754 & 0.9971 \\
PC4 & 0.0115 & 0.0029 & 1.0000 \\
\bottomrule
\end{tabular}
\end{table}

\begin{table}[htbp]
\centering
\caption{PCA loadings; each column is one component, headed by its share of total variance.}\label{tab:pca-loadings}
\begin{tabular}{lrrrr}
\toprule
variable & PC1 (49.9\%) & PC2 (32.3\%) & PC3 (17.5\%) & PC4 (0.3\%) \\
\midrule
temp & 0.705 & -0.002 & 0.068 & -0.706 \\
atemp & 0.705 & -0.034 & 0.046 & 0.707 \\
hum & -0.075 & -0.703 & 0.707 & -0.005 \\
windspeed & -0.039 & 0.71 & 0.702 & 0.027 \\
\bottomrule
\end{tabular}
\end{table}


\begin{figure}[H]
  \centering
  \includegraphics[width=0.65\textwidth]{figures/scree.pdf}
  \caption{Scree plot of the weather-block PCA, with the Kaiser criterion
  ($\lambda = 1$) marked.}
  \label{fig:scree}
\end{figure}

Although the components are not used in the regression, the PCA
reduces the weather block to three interpretable axes that together
account for $99.7\%$ of its variance. PC1 ($49.9\%$) loads almost
entirely on \textit{temp} and \textit{atemp}, at $0.705$ each, with
humidity and windspeed contributing nothing; it is a pure temperature
axis. PC2 ($32.3\%$) sets humidity against windspeed ($-0.703$ and
$0.710$), separating muggy, still days from dry, windy ones. PC3
($17.5\%$) loads the same two variables in the same direction ($0.707$
and $0.702$), and so measures how unsettled the day is overall rather
than which of the two conditions dominates.

PC4 is the informative null result. It accounts for $0.3\%$ of the
variance and is almost exactly the contrast between \textit{temp} and
\textit{atemp} ($-0.706$ against $0.707$), so the only direction in
which the two temperature measures disagree is one the data barely
varies along. This is the same conclusion the VIF reaches in the
preceding section, and it justifies carrying only \textit{temp} into
the model.

\section{SGDRegressor}
\label{sec:sgd}

The data are split $80/20$ into a training and a held-out test set
under a fixed seed, and the model is selected by grid search over the
training half alone. The search ranges over both available penalties,
\texttt{l1} and \texttt{l2}, and over $100$ logarithmically spaced
values of $\alpha$ from $10^{-5}$ to $10^{3}$, scoring each candidate
by five-fold cross-validated $R^2$. The convergence budget is at most $5000$
epochs at a tolerance of $10^{-4}$, this was chosen so that the
differences between grid points reflect the penalty being tested
rather than one candidate stopping earlier than another. The winning
configuration is refit on the full training set and scored once
against the test set. The cross validation grid search utilizes the
$R^2$ of the model.

\begin{table}[htbp]
\centering
\caption{Tuned SGDRegressor: selected hyper-parameters and held-out performance.}\label{tab:sgd-metrics}
\begin{tabular}{lr}
\toprule
statistic & value \\
\midrule
penalty & l1 \\
alpha & 1.000e-05 \\
CV $R^2$ & 0.6788 \\
test $R^2$ & 0.6775 \\
test RMSE & 101.1 \\
test MAE & 74.2 \\
test MSE & 10211.8 \\
\bottomrule
\end{tabular}
\end{table}


\begin{table}[htbp]
\centering
\caption{SGDRegressor coefficients on standardised predictors, by magnitude.}\label{tab:sgd-coefs}
\begin{tabular}{lr@{\hskip 2.5em}lr}
\toprule
predictor & coefficient & predictor & coefficient \\
\midrule
hr\_17 & 75.446 & hr\_22 & 14.041 \\
hr\_18 & 70.672 & hum & -12.9 \\
hr\_8 & 63.453 & season\_3 & 11.998 \\
hr\_19 & 48.472 & weekday\_6 & 7.362 \\
hr\_16 & 45.469 & hr\_4 & -7.111 \\
temp & 45.46 & hr\_6 & 6.852 \\
yr & 43.01 & hr\_3 & -6.673 \\
hr\_13 & 36.029 & hr\_23 & 6.362 \\
hr\_12 & 35.119 & hr\_2 & -5.448 \\
hr\_9 & 33.333 & weekday\_5 & 5.303 \\
hr\_15 & 32.351 & weathersit\_2 & -4.712 \\
hr\_7 & 31.746 & weekday\_2 & 3.932 \\
hr\_20 & 31.364 & weekday\_1 & 3.924 \\
hr\_14 & 30.393 & weekday\_4 & 3.849 \\
season\_4 & 27.981 & holiday & -3.827 \\
hr\_11 & 26.766 & hr\_5 & -3.692 \\
hr\_21 & 22.614 & hr\_1 & -3.639 \\
hr\_10 & 21.188 & windspeed & -3.167 \\
season\_2 & 18.489 & weekday\_3 & 1.781 \\
weathersit\_3 & -17.037 &  &  \\
\bottomrule
\end{tabular}
\end{table}


The search settles on \texttt{l1} at $\alpha = 10^{-5}$, which is the
smallest value the grid offers. This is the
first of several signs that the model is not variance-limited, which
is further addressed in the regularised OLS in Section~\ref{sec:reg}
makes directly. The cross-validated
$R^2$ of $0.6788$ and the test $R^2$ of $0.6775$ agree to within
$0.0013$, so the grid search has not tuned itself into the folds.

The error metrics are reported in riders per hour and are more
informative here than $R^2$. Mean ridership over the window is about
$189$ riders/hour, so an RMSE of $101.1$ is substantial since a typical
prediction misses by more than half the mean. That RMSE exceeds MAE
says that a minority of hours are missed badly
enough to dominate the squared error. This is likely because
ridership depending on the hour depends on season in a non trivial
manner, $8$ in the peak of the summer has much more light than 8
during the winter, where it is pitch black. As such, the model splits
the difference and underpredicts the former and over predicts the latter.

\section{OLS: \texttt{fit}}
\label{sec:ols}

Where \texttt{SGDRegressor} approaches the coefficients iteratively,
this section solves for them in closed form, and does so on all
$17{,}379$ observations rather than the training split, in order to
obtain standard errors, $t$-statistics and confidence intervals for
every predictor. The reported $R^2$ is consequently in-sample and is
not directly comparable with the held-out figures above. That it
exceeds the SGD test $R^2$ by only $0.0044$ shows that almost none of
the fit is specific to the sample it was
estimated on.

\begin{table}[htbp]
\centering
\caption{OLS (\texttt{fit}) overall fit statistics and residual diagnostics.}\label{tab:ols-fitstats}
\begin{tabular}{lr}
\toprule
statistic & value \\
\midrule
observations & 17379 \\
df model & 39 \\
df residuals & 17339 \\
$R^2$ & 0.6819 \\
adjusted $R^2$ & 0.6811 \\
F-statistic & 952.9 \\
Prob (F) & 0.000e+00 \\
log-likelihood & -105089.2 \\
AIC & 210258.5 \\
BIC & 210569.0 \\
Durbin--Watson & 0.4982 \\
Jarque--Bera & 2624.3 \\
Prob (JB) & 0.000e+00 \\
skew & 0.4724 \\
kurtosis & 4.6527 \\
condition number & 6.8 \\
\bottomrule
\end{tabular}
\end{table}


\begin{figure}[H]
  \centering
  \includegraphics[width=0.6\textwidth]{figures/qqplot.pdf}
  \caption{Normal Q--Q plot of the OLS residuals.}
  \label{fig:qq}
\end{figure}

\begin{figure}[H]
  \centering
  \includegraphics[width=\textwidth]{figures/hour_effects.pdf}
  \caption{Observed mean ridership by hour of day (bars, left axis) against the
    fitted OLS hour coefficients (line, right axis). Coefficients are
    rescaled to riders/hour and are relative to the dropped reference
  level, midnight.}
  \label{fig:hour-effects}
\end{figure}

\begin{longtable}{lrrrrrr}
\caption{OLS (\texttt{fit}) coefficient table with standard errors, $t$-statistics, $p$-values and 95\% confidence intervals.}\label{tab:ols-coefs}\\
\toprule
predictor & coef & std err & t & $P>|t|$ & [0.025 & 0.975] \\
\midrule
\endfirsthead
\toprule
predictor & coef & std err & t & $P>|t|$ & [0.025 & 0.975] \\
\midrule
\endhead
\midrule \multicolumn{7}{r}{\footnotesize continued} \\
\endfoot
\bottomrule
\endlastfoot
yr & 42.743 & 0.784 & 54.503 & 0.000e+00 & 41.206 & 44.28 \\
holiday & -4.834 & 0.812 & -5.956 & 2.638e-09 & -6.425 & -3.243 \\
temp & 47.176 & 1.369 & 34.459 & 9.469e-252 & 44.492 & 49.859 \\
hum & -13.272 & 1.049 & -12.648 & 1.654e-36 & -15.329 & -11.216 \\
windspeed & -4.152 & 0.838 & -4.955 & 7.306e-07 & -5.794 & -2.509 \\
season\_2 & 18.658 & 1.228 & 15.194 & 8.357e-52 & 16.251 & 21.065 \\
season\_3 & 11.976 & 1.598 & 7.496 & 6.875e-14 & 8.845 & 15.108 \\
season\_4 & 28.233 & 1.047 & 26.97 & 5.744e-157 & 26.181 & 30.285 \\
hr\_1 & -3.499 & 1.075 & -3.254 & 0.001 & -5.606 & -1.391 \\
hr\_2 & -5.298 & 1.072 & -4.941 & 7.852e-07 & -7.4 & -3.196 \\
hr\_3 & -7.325 & 1.067 & -6.866 & 6.817e-12 & -9.416 & -5.234 \\
hr\_4 & -7.964 & 1.068 & -7.459 & 9.107e-14 & -10.056 & -5.871 \\
hr\_5 & -4.735 & 1.075 & -4.405 & 1.063e-05 & -6.842 & -2.628 \\
hr\_6 & 7.014 & 1.078 & 6.508 & 7.832e-11 & 4.901 & 9.126 \\
hr\_7 & 34.091 & 1.077 & 31.643 & 1.070e-213 & 31.979 & 36.203 \\
hr\_8 & 62.253 & 1.077 & 57.828 & 0.000e+00 & 60.143 & 64.363 \\
hr\_9 & 32.725 & 1.078 & 30.365 & 2.287e-197 & 30.612 & 34.837 \\
hr\_10 & 21.802 & 1.081 & 20.16 & 2.342e-89 & 19.682 & 23.922 \\
hr\_11 & 26.943 & 1.088 & 24.764 & 4.458e-133 & 24.811 & 29.076 \\
hr\_12 & 34.88 & 1.096 & 31.818 & 5.699e-216 & 32.732 & 37.029 \\
hr\_13 & 33.921 & 1.103 & 30.758 & 2.490e-202 & 31.759 & 36.082 \\
hr\_14 & 30.75 & 1.108 & 27.76 & 5.480e-166 & 28.578 & 32.921 \\
hr\_15 & 32.632 & 1.109 & 29.417 & 1.174e-185 & 30.458 & 34.807 \\
hr\_16 & 45.093 & 1.108 & 40.703 & 0.000e+00 & 42.921 & 47.264 \\
hr\_17 & 75.904 & 1.103 & 68.821 & 0.000e+00 & 73.742 & 78.065 \\
hr\_18 & 69.409 & 1.096 & 63.346 & 0.000e+00 & 67.261 & 71.556 \\
hr\_19 & 47.616 & 1.087 & 43.791 & 0.000e+00 & 45.485 & 49.747 \\
hr\_20 & 31.656 & 1.083 & 29.244 & 1.509e-183 & 29.535 & 33.778 \\
hr\_21 & 21.717 & 1.079 & 20.13 & 4.227e-89 & 19.603 & 23.832 \\
hr\_22 & 14.274 & 1.077 & 13.25 & 7.030e-40 & 12.163 & 16.386 \\
hr\_23 & 6.469 & 1.077 & 6.009 & 1.909e-09 & 4.359 & 8.579 \\
weekday\_1 & 3.287 & 1.045 & 3.145 & 0.002 & 1.238 & 5.336 \\
weekday\_2 & 3.7 & 1.016 & 3.643 & 2.706e-04 & 1.709 & 5.692 \\
weekday\_3 & 4.606 & 1.017 & 4.527 & 6.019e-06 & 2.612 & 6.601 \\
weekday\_4 & 4.737 & 1.017 & 4.658 & 3.222e-06 & 2.744 & 6.731 \\
weekday\_5 & 5.94 & 1.018 & 5.837 & 5.396e-09 & 3.946 & 7.935 \\
weekday\_6 & 5.684 & 1.018 & 5.583 & 2.396e-08 & 3.689 & 7.68 \\
weathersit\_2 & -4.769 & 0.847 & -5.63 & 1.825e-08 & -6.43 & -3.109 \\
weathersit\_3 & -18.216 & 0.889 & -20.482 & 3.838e-92 & -19.959 & -16.473 \\
const & 189.463 & 0.777 & 243.857 & 0.000e+00 & 187.94 & 190.986 \\
\end{longtable}


The model as a whole is significant beyond any useful resolution, $F =
952.9$ on $39$ and $17{,}339$ degrees of freedom, with a $p$-value that
underflows to zero. The adjusted $R^2$ of $0.6811$ falls only $0.0008$
below the unadjusted $0.6819$, so the $39$ predictors are well justified.

Every predictor in Table~\ref{tab:ols-coefs} is significant at the
$5\%$ level, and all but two clear $10^{-3}$. The weakest is
\textit{weekday\_1} at $p = 0.002$; the strongest are reported as
exactly zero because the $p$-values fall below double precision. The
confidence intervals are correspondingly narrow. The interval on
\textit{temp} spans $[44.49,\ 49.86]$ around a point estimate of
$47.18$, a width of about $11\%$ of the estimate, and the interval on
\textit{yr} spans $[41.21,\ 44.28]$ around $42.74$. Nothing in the
table is close to admitting zero.

Because the design matrix is standardised, each coefficient is the
change in ridership associated with a one standard deviation change in
that predictor, holding the rest fixed. This is natural for the
continuous weather variables and awkward for the one-hot columns, where
a standard deviation is a fraction of a level; Figure~\ref{fig:hour-effects}
therefore rescales the hour coefficients back into riders per hour
against the midnight reference. On that scale the 5~p.m. effect is
close to $380$ riders/hour above midnight and the 8~a.m. effect close
to $310$. These hours make a great deal of sense as people commuting
to and from work. Hour of day is by a wide margin the dominant predictor, and
the fitted hour profile tracks the observed hourly means closely,
including the twin commute peaks and the midday shoulder between them.
This is the confirmation promised in Section~\ref{sec:features}, the
relationship
between hour and ridership is not monotone, so an ordinal encoding
would have been unable to represent it and one-hot encoding was the
correct choice.

Among the remaining predictors, \textit{temp} is the largest at
$47.18$ riders/hour per standard deviation, and the year indicator is
nearly as large at $42.74$, which is the level shift already visible in
Figure~\ref{fig:timeseries} and the reason \textit{yr} was retained
despite not being of interest in itself. Autumn contributes $28.23$
over the spring reference. The weather terms all suppress ridership in
the expected direction and order: bad weather costs $18.22$, humidity
$13.27$, wind $4.15$, and a holiday $4.83$. Day-of-week effects are
significant but small, spanning only $3.29$ to $5.94$ across the six
non-reference days, so once the hour, the season and the weather are
known, which weekday it is adds little.

The residuals are also non-normal. The Jarque--Bera statistic is
$2624.3$ with skew $0.4724$ and kurtosis $4.6527$, describing a
right-skewed, heavy-tailed distribution, and the Q--Q plot in
Figure~\ref{fig:qq} shows the departure concentrated in the upper tail.
This is the same phenomenon as the RMSE MAE gap in the previous
section, the model under-predicts the
busiest hours leading to heavy tails at the extreme of the residual
distribution, and ridership is a count bounded below at zero whose
variance grows with its mean, which a homoskedastic linear model in the
levels cannot accommodate. Plausibly, a log transform of the response would fit
the shape of the data better. Collinearity, by contrast, is no longer a
concern: the condition number of the fitted standardised design is
$6.8$, against the $3.33 \times 10^{15}$ recorded in
Table~\ref{tab:cond} before the drops.

\section{OLS: \texttt{fit\_regularized}}
\label{sec:reg}

The penalised fits are estimated on the training split only and scored
by $R^2$ on the held-out test set, so that the comparison between
penalty strengths is an out-of-sample one. The grid crosses seven
values of $\alpha$, from $0$ to $10$, with five values of
\texttt{L1\_wt} from pure ridge to pure lasso. The penalty vector is
constructed so that the intercept is left unpenalised, which matters
here: the intercept carries the mean ridership of roughly $189$
riders/hour, and shrinking it toward zero would force the other
coefficients to absorb a large constant.

\begin{table}[htbp]
\centering
\caption{Regularised OLS (\texttt{fit\_regularized}): held-out $R^2$ over the $(\alpha,\ \mathrm{L1\_wt})$ grid. Columns are L1\_wt.}\label{tab:reg-grid}
\begin{tabular}{lrrrrr}
\toprule
$\alpha$ & 0 & 0.15 & 0.5 & 0.85 & 1 \\
\midrule
0 & 0.6783 & 0.6783 & 0.6783 & 0.6783 & 0.6783 \\
0.0001 & 0.6783 & 0.6783 & 0.6783 & 0.6783 & 0.6783 \\
0.001 & 0.6783 & 0.6783 & 0.6783 & 0.6783 & 0.6783 \\
0.01 & 0.6778 & 0.6779 & 0.6782 & 0.6783 & 0.6783 \\
0.1 & 0.6642 & 0.6661 & 0.6641 & 0.6767 & 0.6779 \\
1 & 0.5263 & 0.5468 & 0.5986 & 0.6512 & 0.6563 \\
10 & 0.1447 & 0.1587 & 0.2146 & 0.3799 & 0.5922 \\
\bottomrule
\end{tabular}
\end{table}

\begin{table}[htbp]
\centering
\caption{Regularised OLS: grid point with the highest held-out $R^2$.}\label{tab:reg-best}
\begin{tabular}{lr}
\toprule
statistic & value \\
\midrule
$\alpha$ & 0.001 \\
L1\_wt & 0.15 \\
test $R^2$ & 0.6783 \\
test MSE & 10186.4 \\
test RMSE & 100.9 \\
non-zero coefficients & 40/40 \\
\bottomrule
\end{tabular}
\end{table}


The grid is flat where it is not actively harmful. Held-out $R^2$ is
$0.6783$ at every combination with $\alpha \le 10^{-3}$, and the best
grid point, $\alpha = 10^{-3}$ with $\texttt{L1\_wt} = 0.15$, matches
the wholly unpenalised fit at $\alpha = 0$ to four decimal places. At
the selected point all $40$ coefficients remain non-zero: the lasso
component eliminates nothing. Past $\alpha = 10^{-2}$ the fit degrades
monotonically, and by $\alpha = 10$ ridge has collapsed to $R^2 =
0.1447$. Lasso degrades more gracefully at the top of the range, at
$0.5922$, because it discards the weaker predictors outright while
leaving the few dominant hour and temperature terms near their
unpenalised values, whereas ridge shrinks everything at once, the
terms carrying the most signal included.

The reason regularisation buys nothing here is that neither condition
it corrects is present. Shrinkage trades bias for variance, and is
most helpful when predictors are collinear, when
they are numerous relative to the sample, or when many of them are
noise. Collinearity was removed by hand in
Section~\ref{sec:features}, leaving a maximum
VIF of $4.23$ and a condition number of $38.4$. The sample provides
$17{,}379$ observations for $39$ parameters, a ratio of over $440$ to
one. And, as Section~\ref{sec:ols} establishes, every predictor is individually
significant, so there is no noise column for the lasso to find and
discard.

\section{Conclusion}
\label{sec:conclusion}

The two estimators agree closely with a test $R^2$ of $0.6775$ for the
tuned SGD fit against $0.6783$ for the best penalised OLS and $0.6819$
in-sample for the unpenalised one. They agree on the substance
as well. Roughly two thirds of the variation in hourly ridership is
explained by when in the day it is, what the weather is doing, and
where in the year the day falls, in that order of importance. The hour
profile is the dominant structure and is primarily explained by
commutors forming twin peaks at 8~a.m. and 5~p.m., a midday shoulder,
and near-zero overnight traffic. Temperature is the strongest
continuous driver, with humidity and bad weather suppressing
ridership and wind mattering comparatively little. Day of week, once
those are accounted for, is almost irrelevant.

The remaining third of the variance is not noise, the residuals are
autocorrelated and right-skewed, which point towards a limitation of
additive model that the levels of these variables cannot
let the size of the evening peak depend on the
weather or the season. An interaction between hour and temperature, or
a count model fitted to the log of the response, would be the natural
next step.

\end{document}
